% WIZM Yellow Paper — formal, current production specification.
% Compile: pdflatex yellowpaper.tex (twice for references).
\documentclass[11pt,a4paper]{article}
\usepackage[utf8]{inputenc}
\usepackage[T1]{fontenc}
\usepackage{lmodern}
\usepackage{amsmath,amssymb,amsthm}
\usepackage{booktabs}
\usepackage{array}
\usepackage{longtable}
\usepackage{geometry}
\usepackage{hyperref}
\usepackage{enumitem}
\usepackage{xcolor}
\usepackage{microtype}
\geometry{margin=24mm}
\hypersetup{colorlinks=true,linkcolor=blue!60!black,citecolor=blue!60!black,urlcolor=blue!60!black}

\theoremstyle{definition}
\newtheorem{definition}{Definition}[section]
\newtheorem{protocol}{Protocol}[section]
\theoremstyle{plain}
\newtheorem{proposition}{Proposition}[section]
\newtheorem{lemma}{Lemma}[section]

\newcommand{\B}{\{0,1\}}
\newcommand{\Bclx}{\B^{160}}
\newcommand{\Bccxxxii}{\B^{256}}
\newcommand{\concat}{\mathbin{\Vert}}

\title{WIZM Yellow Paper\\ \large State Transition, PoA Consensus, Fees, and Transaction Layouts}
\author{WIZM Protocol Engineering}
\date{29 October 2026}

\begin{document}
\maketitle

\begin{abstract}
This document specifies the current WIZM chain: a five-validator
proof-of-authority ledger with transparent and stealth transfers. The
chain uses Falcon-512 signatures, ML-KEM for encrypted mempool and
transport key establishment, BLAKE3 hashing, and AES-256-GCM mempool
encryption. It contains no zero-knowledge, shielded, confidential, or
note-based transaction path. Stealth transfers provide recipient
unlinkability only; sender and amount are public.
\end{abstract}

\section{Network parameters}

\begin{center}
\begin{tabular}{@{}ll@{}}
\toprule
Parameter & Value \\
\midrule
Chain ID & 7777 \\
Validators & 5 PoA seats \\
Steady slot & 10 seconds \\
Genesis slot & 240 seconds \\
Slow-start slot & 60 seconds \\
Max transactions per block & 1200 \\
Max block body & 8 MiB \\
Max sealed throughput & 120 TPS \\
Base fee & 10,000 base units \\
Burst-free transactions per wallet/block & 20 \\
Burst multiplier & 10x per excess transaction \\
\bottomrule
\end{tabular}
\end{center}

The proposal cap is the count cap first for transparent and stealth
transactions. A block is invalid if it exceeds 1200 transactions or
8 MiB, or if a wallet exceeds the burst schedule.

\section{Cryptographic functions}

Let $\mathrm{BLAKE3}:\B^*\to\B^{256}$ be BLAKE3-256. An address is
\[
  A(pk) = \mathrm{BLAKE3}(pk)[0:20] \in \Bclx.
\]
Falcon-512 provides signatures:
\[
  \mathrm{Verify}_{pk}(m,\sigma)=1.
\]
Public keys are 897 bytes; signatures are 512--752 bytes. ML-KEM-512 is
used for encrypted mempool/committee key establishment; ML-KEM-768 is
used for transport. AES-256-GCM encrypts mempool payloads.

\section{State}

A state $\sigma$ maps a 20-byte address to
\[
  \sigma(a) = (balance, nonce, public\_key).
\]
It also contains spent stealth key images, block metadata, committee
state, and tokenomics counters. There is no shielded commitment tree,
note set, nullifier set, or Merkle proof state.

For a sender $a$ and amount $v$ with fee $f$, a transparent transfer is
valid only if
\[
  balance(a) \ge v+f, \qquad nonce(tx)=nonce(a).
\]
After execution:
\[
  balance(a) \leftarrow balance(a)-v-f,\quad
  nonce(a)\leftarrow nonce(a)+1.
\]

\section{Transaction layouts}

All integers are little-endian. Falcon signatures are variable-length.

\subsection{Transparent transfer}
\[
[from:20]\concat[to:20]\concat[amount:8]\concat[nonce:8]
\concat[pk:897]\concat[sig:512..752].
\]
Wire size is 1,465--1,705 bytes. Sender, recipient, and amount are
public. The signature covers the first 56 bytes.

\subsection{Stealth transfer}
\[
[sender:20]\concat[stealth:20]\concat[ephemeral:32]
\concat[view\_tag:1]\concat[amount:8]\concat[key\_image:32]
\concat[pk:897]\concat[sig:512..752].
\]
Wire size is 1,522--1,762 bytes. The 113-byte body is signed. Sender
and amount are public. The one-time \texttt{stealth} address is derived
from the recipient's spend key and an ephemeral shared value. The
key image prevents a second spend. There is no encrypted amount and no
zero-knowledge proof.

The address derivation is
\[
  shared = \mathrm{BLAKE3}(ephemeral),
\]
\[
  stealth = \mathrm{BLAKE3}(spend\_pk \concat shared)[0:20],
\]
\[
  view\_tag = shared[0].
\]
The recipient scans with the view secret and recomputes the expected
address. The sender is visible in the transaction.

\section{PoA consensus}

Validators take turns producing blocks in ten-second slots. A producer
drains at most 1200 transactions from the encrypted mempool, seals a
block, and broadcasts it. Validators verify the block and sign a
notarization vote over the chain/genesis domain, block hash, and height.
A majority of active seats advances finality. This is a crash-fault
finality rule, not a Byzantine-fault-tolerant consensus proof.

Competing blocks at the same height are resolved deterministically:
the earlier seal wins, with a deterministic tie-break and a bounded
unwind. Nodes synchronize full block bodies through unicast and gossip;
they update their mesh-best view only from locally validated blocks.

\section{Fees and tokenomics}

Let $n$ be the wallet's index among its fee-paying transactions in a
block. The burst multiplier is
\[
  \mu(n)=
  \begin{cases}
    1, & n \le 20,\\
    10, & n=21,\\
    100, & n=22,\\
    1000, & n=23,\\
    \dots
  \end{cases}
\]
The charged fee is
\[
  f(n)=10{,}000\cdot\mu(n).
\]
During emission, fees are burned. After emission, fees are distributed
to validators. The block reward is
\[
  R_{validator} = \frac{E_{period}}{B_{period}\cdot N_{validators}}
\]
where $E_{period}$ is the emission for the current two-year period,
$B_{period}=6{,}311{,}520$ blocks, and $N_{validators}=5$. At the
current period this is 88,126,450 base units, or 0.8812645 WIZM per
validator per block.

\section{Capacity}

For a block with $T$ transactions and body size $S$:
\[
  T \le 1200, \qquad S \le 8\cdot 1024^2.
\]
At a ten-second slot,
\[
  \Theta_{\max}=\frac{1200}{10}=120\ \mathrm{TPS}.
\]
A generator may submit faster than the chain seals; those transactions
remain in the encrypted mempool. The chain itself never seals more than
120 TPS under the current cap. The burst-free wallet limit is 20
transactions per block, so at least 60 distinct wallets are needed to
fill a 1,200-transaction block without multiplied fees.

\section{What was removed}

The following are not part of the current protocol:

\begin{itemize}[leftmargin=*]
  \item shielded transfers, shielded deposits, and unshields;
  \item notes, commitments, nullifiers, and the shielded Merkle tree;
  \item zero-knowledge proofs and proof verification;
  \item encrypted stealth amounts;
  \item shielded wallet, RPC, CLI, and Tauri surfaces.
\end{itemize}

Stealth is not confidential. It provides one-time recipient addresses
and key-image double-spend protection only. Sender and amount are public
consensus data.

\section{Verification obligations}

A conforming implementation must reject:

\begin{itemize}[leftmargin=*]
  \item any transaction outside the 1,200/8 MiB block limit;
  \item any stealth transaction outside 1,522--1,762 bytes;
  \item any wallet above the 20-transaction burst-free window without
        the multiplied fee;
  \item any replay or duplicate key image;
  \item any attempt to introduce removed shielded/ZK transaction types.
\end{itemize}

The test suite and load harness must exercise only transparent and
stealth transactions.

\end{document}
